\chapter{Numbers, strings and collections}\label{ch:14}
\section*{The working problem}
Learn the actual value model before building conversion APIs.

MetaML Value is a tagged runtime object. Arithmetic, string primitives, tuples and lists have language-level rules. Rendering is for inspection; host conversion should use the supported typed adapters where possible.

\section*{The runtime value boundary}
Value has a kind tag plus scalar fields or an owned object. Strings, lists, aggregates, closures, constructors, modules and foreign objects have different representation classes. Use helpers such as str\_of, list\_to\_vec and agg\_of only after establishing the appropriate kind. A handwritten cast that guesses representation can turn a language error into native undefined behavior.

\section*{In this chapter}
\begin{enumerate}
\item Integer arithmetic
\item Floating-point rendering
\item String primitives operate on bytes
\item Tuples are heterogeneous products
\item Homogeneous lists
\end{enumerate}
Read the workshops in order, then use the chapter review to connect the small examples to an application boundary. Each example is independently runnable; it does not depend on variables introduced in an earlier listing.

\clearpage
\workshop{14.1}{Integer arithmetic}{C++20}
Small integer calculations establish precedence and the numeric representation used by policy code. Keep overflow behavior visible when designing a domain language. MetaML integer arithmetic uses explicit wrapping helpers in the supplied implementation.

\lstinputlisting[language=C++,caption={Example ex14\_01: C++}]{examples/cpp/ex14_01.inc}
\section*{Execution, step by step}
Multiplication binds before addition. The result is an Int value rather than a host double.

\section*{Read the implementation}
The relevant source entry is Session in vendor/MetaML.hpp, around line 8171. Follow the public call into its result construction or dispatch branch. Identify the input ownership, the success condition and the exact failure path before changing this example.

The assertions are the acceptance contract for this workshop. Run the complete program, then change one input at a time. When a change fails, first identify whether the failure belongs to the native build, structural recognition, language checking, execution or the domain result.

\clearpage
\workshop{14.1}{Integer arithmetic}{D}
C++ creates a noncopyable Session value; D allocates MetaML.Session as a class reference. Keep native callback captures alive for the whole session lifecycle and close resources explicitly at their owner boundary. The MetaML program text follows the same language rules in either host. Check the native Result before inspecting the language Value. A successful language value can itself encode a domain failure such as Err or false.

\lstinputlisting[language=D,caption={Example ex14\_01: D}]{examples/d/ex14_01.inc}
\section*{Failure and design boundary}
A mathematically correct identity may be invalid when a domain requires checked rather than wrapping arithmetic.

\section*{Try it yourself}
Compare ordinary addition with Int.checked\_add at the maximum integer boundary.

\section*{Review the result}
Keep the original accepted fixture and add the new boundary fixture beside it. Compare observable values and diagnostic categories across the two hosts. Do not label the D run as passed until a D compiler has executed it; the supplied edition records source review separately from runtime evidence.

The shared language fixture is also available as examples/ex14\_01.ml. Use that file when comparing native embeddings so escaping differences in host strings do not change the program being tested.

\clearpage
\workshop{14.2}{Floating-point rendering}{C++20}
Floating values use a deterministic rendering helper shared by the two dialects. Test boundary values explicitly when output is persisted or compared across hosts. Decimal appearance does not imply decimal arithmetic.

\lstinputlisting[language=C++,caption={Example ex14\_02: C++}]{examples/cpp/ex14_02.inc}
\section*{Execution, step by step}
Both operands are floating literals. The result prints as a shortest round-trip decimal for this finite binary value.

\section*{Read the implementation}
The relevant source entry is Session in vendor/MetaML.hpp, around line 8171. Follow the public call into its result construction or dispatch branch. Identify the input ownership, the success condition and the exact failure path before changing this example.

The assertions are the acceptance contract for this workshop. Run the complete program, then change one input at a time. When a change fails, first identify whether the failure belongs to the native build, structural recognition, language checking, execution or the domain result.

\clearpage
\workshop{14.2}{Floating-point rendering}{D}
C++ creates a noncopyable Session value; D allocates MetaML.Session as a class reference. Keep native callback captures alive for the whole session lifecycle and close resources explicitly at their owner boundary. The MetaML program text follows the same language rules in either host. Check the native Result before inspecting the language Value. A successful language value can itself encode a domain failure such as Err or false.

\lstinputlisting[language=D,caption={Example ex14\_02: D}]{examples/d/ex14_02.inc}
\section*{Failure and design boundary}
Do not use rendered floating strings as a replacement for an application-level tolerance or a specified numeric serialization format.

\section*{Try it yourself}
Compare Float.nan and Float.infinity handling before allowing them in a configuration format.

\section*{Review the result}
Keep the original accepted fixture and add the new boundary fixture beside it. Compare observable values and diagnostic categories across the two hosts. Do not label the D run as passed until a D compiler has executed it; the supplied edition records source review separately from runtime evidence.

The shared language fixture is also available as examples/ex14\_02.ml. Use that file when comparing native embeddings so escaping differences in host strings do not change the program being tested.

\clearpage
\workshop{14.3}{String primitives operate on bytes}{C++20}
Core string primitives provide byte-based operations. ASCII examples make the baseline obvious, while multilingual input requires an explicit UTF-8 policy. Character count and encoded-byte count answer different questions.

\lstinputlisting[language=C++,caption={Example ex14\_03: C++}]{examples/cpp/ex14_03.inc}
\section*{Execution, step by step}
The primitive returns the byte length as an Int. The host adapter sees an ordinary language value.

\section*{Read the implementation}
The relevant source entry is Session in vendor/MetaML.hpp, around line 8171. Follow the public call into its result construction or dispatch branch. Identify the input ownership, the success condition and the exact failure path before changing this example.

The assertions are the acceptance contract for this workshop. Run the complete program, then change one input at a time. When a change fails, first identify whether the failure belongs to the native build, structural recognition, language checking, execution or the domain result.

\clearpage
\workshop{14.3}{String primitives operate on bytes}{D}
C++ creates a noncopyable Session value; D allocates MetaML.Session as a class reference. Keep native callback captures alive for the whole session lifecycle and close resources explicitly at their owner boundary. The MetaML program text follows the same language rules in either host. Check the native Result before inspecting the language Value. A successful language value can itself encode a domain failure such as Err or false.

\lstinputlisting[language=D,caption={Example ex14\_03: D}]{examples/d/ex14_03.inc}
\section*{Failure and design boundary}
The separately supplied metalib may offer higher-level Unicode helpers, but those files are not present in this source set.

\section*{Try it yourself}
Measure a non-ASCII UTF-8 string and explain the result in bytes.

\section*{Review the result}
Keep the original accepted fixture and add the new boundary fixture beside it. Compare observable values and diagnostic categories across the two hosts. Do not label the D run as passed until a D compiler has executed it; the supplied edition records source review separately from runtime evidence.

The shared language fixture is also available as examples/ex14\_03.ml. Use that file when comparing native embeddings so escaping differences in host strings do not change the program being tested.

\clearpage
\workshop{14.4}{Tuples are heterogeneous products}{C++20}
A tuple combines values positionally and can contain different types. It is useful for small internal return values, but a record usually documents a public configuration object more clearly than positional fields.

\lstinputlisting[language=C++,caption={Example ex14\_04: C++}]{examples/cpp/ex14_04.inc}
\section*{Execution, step by step}
The three positions retain their different types. Rendering exposes the product shape without flattening it to a string list.

\section*{Read the implementation}
The relevant source entry is Session in vendor/MetaML.hpp, around line 8171. Follow the public call into its result construction or dispatch branch. Identify the input ownership, the success condition and the exact failure path before changing this example.

The assertions are the acceptance contract for this workshop. Run the complete program, then change one input at a time. When a change fails, first identify whether the failure belongs to the native build, structural recognition, language checking, execution or the domain result.

\clearpage
\workshop{14.4}{Tuples are heterogeneous products}{D}
C++ creates a noncopyable Session value; D allocates MetaML.Session as a class reference. Keep native callback captures alive for the whole session lifecycle and close resources explicitly at their owner boundary. The MetaML program text follows the same language rules in either host. Check the native Result before inspecting the language Value. A successful language value can itself encode a domain failure such as Err or false.

\lstinputlisting[language=D,caption={Example ex14\_04: D}]{examples/d/ex14_04.inc}
\section*{Failure and design boundary}
Host pair conversion covers a specific two-element mapping; do not assume every tuple shape has an automatic C++ struct adapter.

\section*{Try it yourself}
Write a pattern that extracts the second component and verify its inferred type.

\section*{Review the result}
Keep the original accepted fixture and add the new boundary fixture beside it. Compare observable values and diagnostic categories across the two hosts. Do not label the D run as passed until a D compiler has executed it; the supplied edition records source review separately from runtime evidence.

The shared language fixture is also available as examples/ex14\_04.ml. Use that file when comparing native embeddings so escaping differences in host strings do not change the program being tested.

\clearpage
\workshop{14.5}{Homogeneous lists}{C++20}
A list contains values of one compatible element type. The checker rejects heterogeneous lists rather than inserting implicit dynamic wrappers. Empty lists are polymorphic until context constrains their element type.

\lstinputlisting[language=C++,caption={Example ex14\_05: C++}]{examples/cpp/ex14_05.inc}
\section*{Execution, step by step}
The three integers form a linked runtime list. The pretty printer uses bracket notation.

\section*{Read the implementation}
The relevant source entry is Session in vendor/MetaML.hpp, around line 8171. Follow the public call into its result construction or dispatch branch. Identify the input ownership, the success condition and the exact failure path before changing this example.

The assertions are the acceptance contract for this workshop. Run the complete program, then change one input at a time. When a change fails, first identify whether the failure belongs to the native build, structural recognition, language checking, execution or the domain result.

\clearpage
\workshop{14.5}{Homogeneous lists}{D}
C++ creates a noncopyable Session value; D allocates MetaML.Session as a class reference. Keep native callback captures alive for the whole session lifecycle and close resources explicitly at their owner boundary. The MetaML program text follows the same language rules in either host. Check the native Result before inspecting the language Value. A successful language value can itself encode a domain failure such as Err or false.

\lstinputlisting[language=D,caption={Example ex14\_05: D}]{examples/d/ex14_05.inc}
\section*{Failure and design boundary}
A list value is not a random-access host vector. Convert at the boundary when your native algorithm needs contiguous storage.

\section*{Try it yourself}
Ask type\_of for an empty list and then for a list containing both Int and Bool.

\section*{Review the result}
Keep the original accepted fixture and add the new boundary fixture beside it. Compare observable values and diagnostic categories across the two hosts. Do not label the D run as passed until a D compiler has executed it; the supplied edition records source review separately from runtime evidence.

The shared language fixture is also available as examples/ex14\_05.ml. Use that file when comparing native embeddings so escaping differences in host strings do not change the program being tested.

\clearpage
\section*{Chapter review: The runtime value boundary}
Core integer helpers intentionally implement a wrapping model, while Int.checked\_add, checked\_sub and checked\_mul return Option for checked operations. Floating values need special handling for NaN, infinities and negative zero. String length and substring primitives use byte positions. Host::Conv supplies a documented set of conversions, not arbitrary C++ reflection. Treat Value pass-through as a low-level escape hatch requiring a clear ownership and type contract.

\section*{Questions to answer in your implementation}
\begin{enumerate}
\item Compare ordinary addition with Int.checked\_add at the maximum integer boundary.
\item Compare Float.nan and Float.infinity handling before allowing them in a configuration format.
\item Measure a non-ASCII UTF-8 string and explain the result in bytes.
\item Write a pattern that extracts the second component and verify its inferred type.
\item Ask type\_of for an empty list and then for a list containing both Int and Bool.
\end{enumerate}
\section*{A practical completion test}
Choose one workshop and replace its fixed input with an application boundary: a command-line argument, a file, a configuration field or a typed callback. Write down the accepted input shape and the error outcome. Keep parsing and validation separate from publishing state. Preserve the smallest failing input when a regression is found.

